% Part: computability
% Chapter: recursive-functions
% Section: examples

\documentclass[../../../include/open-logic-section]{subfiles}

\begin{document}

\olfileid{cmp}{rec}{exa}
\olsection{ఆదిమ పునరావృత్త ప్రమేయాల ఉదాహరణలు}


ఆదిమ పునరావృత్త ప్రమేయాలకు కొన్ని ఉదాహరణలు ఇప్పటికే ఉన్నాయి:
కూడిక, గుణకార ప్రమేయాలు~$\Add$, $\Mult$. తాదాత్మ్య ప్రమేయం
$\fn{id}(x)=x$ కేవలం~$\Proj{1}{0}$ కనుక అది ఆదిమ పునరావృత్త
ప్రమేయం. స్థిర ప్రమేయాలు $\fn{const}_n(x)=n$ కూడా ఆదిమ పునరావృత్త
ప్రమేయాలే; ఎందుకంటే $\Zero$, $\Succ$ల వరుస సంయుక్తాలతో వాటిని
నిర్వచించవచ్చు. ఆదిమ పునరావృత్త నిర్వచనాల్లో స్థిరాంకాలను వాడాలంటే
ఇది ఉపయోగపడుతుంది. ఉదాహరణకు $f(x)=2\cdot x$ను నిర్వచించాలంటే,
$\fn{const}_2(x)$, గుణకారం నుంచి సంయుక్తం ద్వారా
$f(x)=\Mult(\fn{const}_2(x),\Proj{1}{0}(x))$గా పొందవచ్చు. ఇకపై ఈ
పద్ధతిని ఉపయోగిస్తాం.
\sourcecorrection{OLTECMPREC-006}{ఆంగ్ల మూలం $f(x)=2\cdot x$ నిర్మాణానికి సాధారణ $\fn{const}_n$ను పేర్కొంది; వెంటనే ఇచ్చిన సమీకరణకు అనుగుణంగా అవసరమైన $\fn{const}_2$గా సరిచేశాం.}

\begin{prop}
  ఘాతాంక ప్రమేయం $\fn{exp}(x,y)=x^y$ ఆదిమ పునరావృత్త ప్రమేయం.
\end{prop}

\begin{proof}
  $\fn{exp}$ను ఆదిమ పునరావృత్తంగా ఇలా నిర్వచించవచ్చు:
  \begin{align*}
    \fn{exp}(x, 0) & = 1\\
    \fn{exp}(x, y+1) & = \Mult(x, \fn{exp}(x,y)).
    \intertext{ఖచ్చితంగా చెప్పాలంటే, ఇది ఆదిమ పునరావృత్త ప్రమేయాల
      నుంచి ఇచ్చిన పునరావృత్త నిర్వచనం కాదు. నియమబద్ధమైన రూపం ఇలా
      ఉంటుంది:}
    \fn{exp}(x, 0) & = f(x)\\
    \fn{exp}(x, y+1) & = g(x, y, \fn{exp}(x,y)).
    \intertext{ఇక్కడ}
    f(x) & = \Succ(\Zero(x)) = 1\\
    g(x, y, z) & = \Mult(\Proj{3}{0}(x, y, z), \Proj{3}{2}(x, y, z)) = x \cdot z
  \end{align*}
  కనుక $f$, $g$లను ఆదిమ పునరావృత్త ప్రమేయాల నుంచి సంయుక్తం ద్వారా
  నిర్వచించాం.
\end{proof}

\begin{prop}
  కింది విధంగా నిర్వచించిన పూర్వగామి ప్రమేయం $\fn{pred}(y)$ ఆదిమ
  పునరావృత్త ప్రమేయం:
  \[
  \fn{pred}(y) = \begin{cases}
    0 & \text{$y=0$ అయితే}\\
    y-1 & \text{ఇతర సందర్భాల్లో.}
  \end{cases}
  \]
\end{prop}

\begin{proof}
 గమనించండి:
 \begin{align*}
   \fn{pred}(0) & = 0 \text{ మరియు}\\
   \fn{pred}(y+1) & = y.
 \end{align*}
 ఇది దాదాపు ఆదిమ పునరావృత్త నిర్వచనమే. అయితే ఖచ్చితంగా చూస్తే,
 ఆదిమ పునరావృత్త నిర్వచన నమూనా కనీసం ఒక అదనపు ఆర్గుమెంటు~$x$ను
 కోరుతుంది కనుక ఇది ఆ నమూనాలో ఇమడదు. $\fn{pred}(y+1)$ నిర్వచనంలో
 $\fn{pred}(y)$ను అసలు వాడకపోవడం కూడా విచిత్రమే. కానీ మొదట
 $\fn{pred}'(x,y)$ను ఇలా నిర్వచించవచ్చు:
 \begin{align*}
   \fn{pred}'(x, 0) & = \Zero(x) = 0,\\
   \fn{pred}'(x, y+1) & = \Proj{3}{1}(x, y, \fn{pred'}(x, y)) = y.
 \end{align*}
తరువాత దాని నుంచి సంయుక్తం ద్వారా, ఉదాహరణకు
$\fn{pred}(x)=\fn{pred}'(\Zero(x),\Proj{1}{0}(x))$గా $\fn{pred}$ను
నిర్వచించవచ్చు.
\end{proof}

\begin{prop}
  క్రమగుణిత ప్రమేయం $\fn{fac}(x)=\fact{x}=1\cdot2\cdot3
  \cdot\dots\cdot x$ ఆదిమ పునరావృత్త ప్రమేయం.
\end{prop}

\begin{proof}
  సహజమైన ఆదిమ పునరావృత్త నిర్వచనం ఇది:
  \begin{align*}
    \fn{fac}(0) &= 1\\
    \fn{fac}(y+1) & = \fn{fac}(y) \cdot (y+1).
    \intertext{నియమబద్ధంగా, మొదట ద్విస్థానిక ప్రమేయం $h$ను
      నిర్వచించాలి:}
    h(x, 0) & = \fn{const}_1(x)\\
    h(x, y+1) & = g(x, y, h(x, y))
    \intertext{ఇక్కడ $g(x,y,z)=\Mult(\Proj{3}{2}(x,y,z),
      \Succ(\Proj{3}{1}(x,y,z)))$గా తీసుకొని, తరువాత}
    \fn{fac}(y) & = h(\Proj{1}{0}(y), \Proj{1}{0}(y)) = h(y,y).
  \end{align*}
  అని తీసుకోవాలి. ఇకపై ఇంత నియమబద్ధమైన సంయుక్త, ఆదిమ పునరావృత్త
  వివరాలన్నిటినీ ప్రతిసారీ ఇవ్వం.
\end{proof}

\begin{prop}
  కింది విధంగా నిర్వచించిన కత్తిరించిన వ్యవకలనం $x\tsub y$ ఆదిమ
  పునరావృత్త ప్రమేయం:
  \[
  x \tsub y = \begin{cases}
    0 & \text{$x<y$ అయితే}\\
    x-y & \text{ఇతర సందర్భాల్లో.}
  \end{cases}
  \]
\end{prop}

\begin{proof}
  మనకు:
  \begin{align*}
    x \tsub 0 & = x\\
    x \tsub (y+1) & = \fn{pred}(x \tsub y)
  \end{align*}
\end{proof}

\begin{prop}
 $x$, $y$ల మధ్య దూరం $\left|x-y\right|$ ఆదిమ పునరావృత్త ప్రమేయం.
\end{prop}

\begin{proof}
  $\left|x-y\right|=(x\tsub y)+(y\tsub x)$. కనుక ఆదిమ పునరావృత్త
  ప్రమేయాలైన $+$, $\tsub$ల నుంచి సంయుక్తం ద్వారా దూరాన్ని
  నిర్వచించవచ్చు.
\end{proof}

\begin{prop}
  $x$, $y$ల గరిష్ఠం $\fn{max}(x,y)$ ఆదిమ పునరావృత్త ప్రమేయం.
\end{prop}

\begin{proof}
  $+$, $\tsub$ల నుంచి సంయుక్తం ద్వారా $\fn{max}(x,y)$ను
  \[
\fn{max}(x,y) \defis x + (y \tsub x).
  \]
  అని నిర్వచించవచ్చు. $x$ గరిష్ఠమైతే, అంటే $x\ge y$ అయితే,
  $y\tsub x=0$; కనుక $x+(y\tsub x)=x+0=x$. $y$ గరిష్ఠమైతే,
  $y\tsub x=y-x$; కనుక $x+(y\tsub x)=x+(y-x)=y$.
\end{proof}

\begin{prop}
  \ollabel{prop:min-pr}
  $x$, $y$ల కనిష్ఠం $\fn{min}(x,y)$ ఆదిమ పునరావృత్త ప్రమేయం.
\end{prop}

\begin{proof}
  అభ్యాసం.
\end{proof}

\begin{prob}
  \olref[cmp][rec][exa]{prop:min-pr}ను నిరూపించండి.
\end{prob}

\begin{prob}
\[
f(x, y) =
2^{(2^{\iddots^{2^{x}}})}\raisebox{1ex}{\bigg\rbrace}
\raisebox{1ex}{\text{$y$ సంఖ్యలో $2$లు}}
\]
ప్రమేయం ఆదిమ పునరావృత్త ప్రమేయమని చూపండి.
\end{prob}

\begin{prob}
పూర్ణాంక భాగహార ప్రమేయం $d(x,y)=\lfloor x/y\rfloor$ (అంటే దశాంశ
బిందువు తరువాతి భాగమంతా విస్మరించే భాగహారం) ఆదిమ పునరావృత్త
ప్రమేయమని చూపండి. $y=0$ అయినప్పుడు $d(x,y)=0$గా నిర్దేశిస్తాం.
ఆదిమ పునరావృత్తి, సంయుక్తం ఉపయోగించి $d$కు స్పష్టమైన నిర్వచనం ఇవ్వండి.
\end{prob}

\begin{prop}
ఆదిమ పునరావృత్త ప్రమేయాల సమితి కింది రెండు క్రియల కింద సంవృతం:
\begin{enumerate}
\item పరిమిత మొత్తాలు: $f(\vec x,z)$ ఆదిమ పునరావృత్త ప్రమేయమైతే,
కింది ప్రమేయం కూడా అలాంటిదే:
\[
g(\vec x, y) \defis \sum_{z = 0}^y f(\vec x, z).
\]
\item పరిమిత లబ్ధాలు: $f(\vec x,z)$ ఆదిమ పునరావృత్త ప్రమేయమైతే,
కింది ప్రమేయం కూడా అలాంటిదే:
\[
h(\vec x, y) \defis \prod_{z = 0}^y f(\vec x, z).
\]
\end{enumerate}
\end{prop}

\begin{proof}
ఉదాహరణకు, పరిమిత మొత్తాలను కింది సమీకరణాలతో పునరావృత్తంగా
నిర్వచించవచ్చు:
\begin{align*}
  g(\vec x, 0) & = f(\vec x, 0)\\
  g(\vec x, y+1) & = g(\vec x, y) + f(\vec x, y+1).
\end{align*}
\end{proof}


\end{document}
